\section{Conclusions} \label{sec:conclusion}
This work presented a two-stage multimodal framework that connects semantic perception with global candidate retrieval. The staged optimization and evaluator-specific reporting distinguish range-view segmentation, native point-level refinement, and external 3D audit results. The segmentation stage preserves useful evidence from persistent and landmark classes, but neither the complete multimodal FFAM++ configuration nor refinement establishes an overall mIoU advantage over the strongest segmentation baselines. Because RGB and FFAM++ are introduced together in the available comparison, the specific contribution of the additional channel gate remains unquantified.

SGMD is not intended to recognize absolute locations or make final loop-closure decisions. On the unseen SemanticKITTI sequence 08, temporal submaps improve spatial Recall@K over individual frames. However, Scan Context and M2DP remain substantially stronger at strict spatial thresholds. SGMD instead produces rankings with high semantic and landmark agreement. Landmark-aware re-ranking increases this agreement and improves strict 5\,m and 10\,m Recall@20 on sequence 08, although the gain is not consistent at the relaxed 20\,m threshold.

The supported contribution is therefore a semantic candidate-retrieval representation that preserves scene composition and persistent landmark evidence. Its role is complementary to geometry-based descriptors rather than a replacement for strict spatial retrieval. Evaluation on cross-session data with seasonal, illumination, or structural changes remains necessary before claiming long-term place-recognition robustness.
